\chapter{Design Documentation --- High Level}
\label{chap:design-high}

\section{System Overview}

The Clue! simulation is a desktop game implemented in Unity (C\#). The table below identifies the major subsystems, their key classes, and responsibilities.

\begin{table}[H]
\centering
\footnotesize
\setlength{\tabcolsep}{5pt}
\begin{tabular}{|p{2.5cm}|p{2.8cm}|p{5.5cm}|}
\hline
\textbf{Subsystem} & \textbf{Key Classes} & \textbf{Responsibility} \\ \hline
Scene Management & \texttt{SceneController}, \texttt{MenuController} & Persists player/AI count between menu and game scenes \\ \hline
Game State & \texttt{GameState} & Tracks Playing / MainMenu / GameOver; controls timescale \\ \hline
Board & \texttt{BoardScript} & 25-row $\times$ 24-col tile grid; single source of truth for occupancy and walkability \\ \hline
Card System & \texttt{Card}, \texttt{Deck} & Loads cards from JSON; creates murder envelope; deals hands \\ \hline
Character Data & \texttt{CharacterData} & ScriptableObject storing character metadata (name, ID) \\ \hline
Game Orchestration & \texttt{GameManagerScript} & Spawns tokens; runs \texttt{MainLoop} coroutine; resolves suggestions and accusations \\ \hline
Player Turn (Human) & \texttt{HumanPlayer} & Waits for UI input via UiManager callbacks \\ \hline
Player Turn (AI) & \texttt{AiPlayer} & Uses BFS-directed movement and knowledge-filtered suggestions; accuses when envelope candidates narrow to 3 \\ \hline
Player Base & \texttt{Player} & Abstract base: turn sequencing, knowledge tracking, notepad state \\ \hline
Movement & \texttt{TokenMover} & Grid movement, room entry/exit, world-space positioning \\ \hline
Rooms / Passages & \texttt{RoomRules} & Static class: room tile detection, secret passage map \\ \hline
Dice & \texttt{DiceRollScript} & Physics Rigidbody dice; result by face orientation \\ \hline
UI & \texttt{UiManager} & Central singleton: all player-facing prompts, suggestion/accusation panels, card reveal \\ \hline
Notebook & \texttt{NotepadGrid} & 147-cell toggle grid for detective deduction \\ \hline
Card Knowledge & \texttt{PlayerKnowledge} & Per-player card inference (Held / NotHeld / disjunctions) \\ \hline
End Screen & \texttt{GameOverScreen} & Restart / return-to-menu on game conclusion \\ \hline
\end{tabular}
\end{table}

\section{Use Cases}

Actors: Human Player, AI Player.

\begin{table}[H]
\centering
\small
\begin{tabular}{|l|l|l|}
\hline
\textbf{Use Case} & \textbf{Human Player} & \textbf{AI Player} \\ \hline
Configure game & \checkmark & --- \\ \hline
Roll dice & \checkmark & \checkmark (auto) \\ \hline
Move token & \checkmark (arrows) & \checkmark (BFS) \\ \hline
Use secret passage & \checkmark & \checkmark (if to target) \\ \hline
Enter room & \checkmark & \checkmark \\ \hline
Make suggestion & \checkmark & \checkmark (filtered) \\ \hline
Respond to suggestion & \checkmark & \checkmark (first card) \\ \hline
View hand & \checkmark & --- \\ \hline
Open detective notebook & \checkmark & --- \\ \hline
Make accusation & \checkmark & \checkmark (3 left) \\ \hline
Restart / return to menu & \checkmark & --- \\ \hline
\end{tabular}
\end{table}

\noindent\textit{Note: ``Make Suggestion'' \texttt{<<include>>} ``Respond to Suggestion''; ``Make Accusation'' ends the game (correct) or eliminates the player (incorrect).}

\begin{figure}[H]
\centering
\includegraphics[width=0.85\textwidth]{use_case_diagram.png}
\caption{Use case diagram --- Human Player and AI Player}
\label{fig:use-case}
\end{figure}

\section{Component Interaction Diagram}

\begin{figure}[H]
\centering
\includegraphics[width=\textwidth]{component_diagram.png}
\caption{Component interaction diagram}
\label{fig:component-diagram}
\end{figure}

\section{Key Design Decisions}

\subsection*{4.1 Unity + C\# (Technology Choice)}
Agreed unanimously at Meeting 1. Unity provides built-in 3D rendering, Rigidbody physics (used for the dice), and cross-platform Mac/PC builds. The team had limited prior Unity experience and saw it as a valuable learning opportunity.

\subsection*{4.2 Singleton Pattern (Manager Classes)}
\texttt{GameManagerScript}, \texttt{BoardScript}, \texttt{UiManager}, \texttt{GameState}, and \texttt{SceneController} all use the Unity Singleton pattern. These represent unique global state that any script may need to access. The singleton prevents duplicate instances and provides a predictable, globally accessible reference.

\subsection*{4.3 Data-Driven Card System}
All card names (characters, weapons, and rooms) are loaded from \texttt{StreamingAssets/cards.json} via \texttt{JsonUtility}. Watson Games explicitly require the game be easily customised. Changing characters or weapons requires only editing this JSON file, with no code changes needed.

\subsection*{4.4 Physics-Based Dice}
\texttt{DiceRollScript} uses a Rigidbody simulation rather than \texttt{Random.Range}. Watson Games require a fun to play and colourful interface, so a physical dice provides visual authenticity. Face orientation is detected by dot product against \texttt{Vector3.up}, and a \texttt{Random.Range(1,7)} fallback handles any degenerate result.

\subsection*{4.5 Callback-Driven UI (Action Delegates)}
\texttt{UiManager} communicates results to game logic via \texttt{System.Action} callbacks, while \texttt{Player} coroutines wait via \texttt{WaitUntil}. This approach keeps UI code and game logic independent.

\subsection*{4.6 AI Player: Basic Decision-Making with Knowledge Scaffold}
\texttt{AiPlayer} uses BFS-directed movement toward envelope-candidate rooms and knowledge-filtered suggestion generation. \texttt{ChooseAccusationFollowUp} returns true when the envelope has been narrowed to exactly 3 candidates. A \texttt{PlayerKnowledge} instance is maintained for each AI player and updated on every suggestion round, providing a foundation for further deduction-based improvements without requiring a rewrite.

\section{Activity Diagram: Full Game Flow}

\begin{figure}[H]
\centering
\includegraphics[width=0.75\textwidth]{activity_game_flow.png}
\caption{Activity diagram --- full game flow}
\label{fig:activity-game-flow}
\end{figure}

\newpage
